\documentclass[12pt]{article}
\usepackage[utf8]{inputenc}
\usepackage[T1]{fontenc}
\usepackage{hyperref}
\title{Isthmus Zapotec Language Resources for K-12 Curriculum Development}
\author{Compiled by Claude Sonnet 5.5 (claude-sonnet-5-5) for Andr\'as Kornai}
\date{}
\begin{document}
\maketitle

\section{Introduction}
Isthmus Zapotec (ISO 639-3 zai, Glottocode isth1244; endonym diidxaz\'a, also called Juchit\'an Zapotec) is an Otomanguean language of the Isthmus of Tehuantepec, Oaxaca, Mexico. Speaker estimates range from about 85,000 (1990 census, per Wikipedia) to about 104,000 (Omniglot). A Latin-alphabet writing system was devised in 1956. The reference description lists it as not known to be taught in schools.

\section{Primary Educational Resources}
The searches found no primers, textbooks or school programmes for Isthmus Zapotec.

\section{Linguistic and Reference Materials}
\begin{itemize}
\item A. Scott Britton, \emph{Zapotec-English/English-Zapotec (Isthmus) Concise Dictionary} (Hippocrene Books, 2003). It has over 5,000 entries and sections on spelling, pronunciation and basic grammar. Listing: \url{https://www.alibri.es/en/libro/568501/zapotec-english-english-zapotec-isthmus-concise-dictionary}
\item Library catalogue record for the same dictionary: \url{https://quaibranly.fr/en/explore-collections/base/Bibliography/action/show/notice/000134836-zapotec-english-english-zapotec-isthmus-concise-dictionary}
\item \emph{Dictionary of the Zapotec Language} in the AILLA archive: a multi-volume Zapotec-Spanish dictionary compiled with collaborators from many communities. It covers several Zapotec varieties, not only Isthmus Zapotec. \url{https://islandora-ailla.lib.utexas.edu/islandora/object/ailla%3A279600}
\item Omniglot, ``Isthmus Zapotec (Diidxaz\'a)'': writing system and background. \url{https://www.omniglot.com/writing/zapotec.php}
\item Wikipedia overview with ISO code data. \url{https://www.wikipedia.com/wiki/ISO_639:zai}
\end{itemize}

\section{Religious and Community Resources}
None found in these searches.

\section{Digital and Audio Resources}
\begin{itemize}
\item El Universal report on a published Zapotec dictionary from a professor at the University of the Tehuantepec Isthmus. The article does not clearly state which Zapotec variety it covers. \url{https://eluniversal.com.mx/english/first-zapotec-dictionary-published}
\end{itemize}

\section{Curriculum Development Context}
The language has a sizeable speaker community and some living use in every generation, but the searches found little classroom material. A curriculum would need new primers and readers. The existing dictionaries and the 1956 orthography are a starting point for the writing system. Resellers listed Britton's dictionary as unavailable or on request, so its availability should be checked.
\end{document}